% !TeX root = statnicaEKT.tex
% !TeX spellcheck = cs_CZ
\part{BPC-MKO - Mobilní komunikace}
\begin{huge}
	Otázky:
	\newline
\end{huge}
\begin{enumerate}
	\item Otázka - Blokové schéma digitálního radiokomunikačního systému, popis bloků, přenosová kapacita systému.
	\item Otázka - Základní způsoby zpracování signálů v systémech mobilních komunikací (zdrojové kódování, kanálové kódování, prokládání a digitální modulace)
	\item Otázka - Systémy s mnohonásobným přístupem a metody multiplexování, proces rozprostírání signálu, výhody a nevýhody jednotlivých přístupů
	\item Otázka - Způsoby přenosu (simplex, duplex) a typy spojování (komutované a paketové), diverzitní příjem, diverzitní systémy
	\item Otázka - Struktura obecné buňkové sítě a její hlavní výhody, procedura "handover" a její různé typy
	\item Otázka - Základní architektura a vlastnosti systémů GSM, GPRS, EDGE, UMTS
	\item Otázka - Systém LTE, základní architektura a popis systému LTE, rozdíly mezi LTE a LTE-A
	\item Otázka - Systém LTE, základní popis E-UTRA Downlink a Uplink a rodzíly mezi nimi, základní popis systémů LTE-CAT, NB-IoT
	\item Otázka - Systémy WLAN (základní popis standardu IEEE 802.11 a definice fyzické vrstvy) a WPAN (základní popis, systém Bluetooth a topologie sítě Bluetooth)
	\item Otázka - Systém LPWAN (základní popis, IoT, systémy LoRa a SigFox)
\end{enumerate}
\newpage
\chapter{První otázka}
\section{Blokové schéma digitálního radiokomunikačního systému}

\begin{figure}[H]
	\centering
	\includegraphics[width=0.7\linewidth]{"Obrazky/MKO/1. otazka/radiokomSystem"}
	\caption{Blokové schéma digitálního radiokomunikačního systému}
	\label{fig:radiokomsystem}
\end{figure}
\begin{itemize}
	\item \textbf{Zdroj signálu} - převede informaci na elektrický signál
	\begin{itemize}
		\item např. mikrofon
	\end{itemize}
	\item \textbf{Převod A/D} (pokud je vstupní signál analogový)
	\item \textbf{Kodér zdroje}
	\begin{itemize}
		\item Kompresí potlačuje redundandní (nadbytečné) a irelevantní (zbytečné) složky v signálu
		\item Použití při přenosu hovorových a přenosových signálů $\rightarrow$ nižší bitová rychlost
	\end{itemize}
	\item \textbf{Kodér kanálu}
	\begin{itemize}
		\item potlačuje chyby při přenosu zvýšením redundance
		\item součástí bývá prokládání (viz KS1) - zabezpečení signálu proti shluku chyb
	\end{itemize}
	\item \textbf{Modulátor}
	\begin{itemize}
		\item modulace na vysokofrekvenční nosnou vlnu (amplitudové ASK, fázové PSK, frekvenční FSK klíčování, M-QAM modulace)
	\end{itemize}
	\item \textbf{Rádiové prostředí} - prostředí pro přenos
	\begin{itemize}
		\item k signálu se cestou přičítá šum, vzniká mezisymbolová interference ISI (viz KS1)
	\end{itemize}
	\item \textbf{Demodulátor}
	\begin{itemize}
		\item převedení z VF zpět ndo základního pásma
		\item extrahování dat původního signálu z nosné vlny
	\end{itemize}
	\item \textbf{Dekodér kanálu}
	\begin{itemize}
		\item detekuje a opravuje chyby (způsobené průchodem radiového prostředí)
	\end{itemize}
	\item \textbf{Dekodér zdroje}
	\begin{itemize}
		\item dekomprese dat diferenčním DPCM
		\item Hybridní kódování (predikce+transformace+kvantování+entropické kódování)
	\end{itemize}
	\item \textbf{D/A převod}
	\begin{itemize}
		\item rekonstrukce z diskrétních vzorků do spojité podoby
	\end{itemize}
	\item \textbf{Příjem - Koncový stupeň}
	\begin{itemize}
		\item např. reporoduktor
	\end{itemize}
\end{itemize}
\section{Přenosová kapacita systému}
Maximální teoretická rychlost, jakou lze bezchybně přenášet data kanálem za přítomnosti šumu.
Závisí na:
\begin{itemize}
	\item přidělené šířce pásma kanálu (B)
	\item poměru výkonu signálu ku šumu S/N
\end{itemize}
\begin{equation}
	C = B \cdot \log_2\left(1 + \frac{S}{N}\right)\,[\mathrm{bit/s}]
\end{equation}

\chapter{Druhá otázka}
\section{Základní způsoby zpracování signálů v systémech mobilních komunikací}
\subsection{Zdrojové kódování - Komprese}
Provádíme za účelem odstranění (u řeči, zvuku, obrazu, videu):
\begin{itemize}
	\item Redundance (nadbytečná data) - data, co se opakují/lze je odvodit od ostatních; vratný proces, \underline{bezztrátová komprese}
	\item Irelevance (zbytečná data) - data, která lidské smysly nedokáží vnímat, nebo nejsou důležité; nevratný proces \underline{ztrátová komprese}
\end{itemize}
Výsledkem je snížení přenosové rychlosti signálu (měnší šířka pásma přenos. k.)
Např.: MPEG, MP3, JPEG, PCM/DPCM\\

Kompresní poměr \textbf{CR} - pro posouzení zdrojového kódování
\begin{equation}
	CR = \frac{R_{\mathrm{vst}}}{R_{\mathrm{vyst}}}\,[-]
\end{equation}
\[R_{\mathrm{vst}} \text{ je přenosová rychlost na vstupu kóderu zdroje } [\mathrm{bit/s}],\]
\[R_{\mathrm{vyst}} \text{ je přenosová rychlost na výstupu kóderu zdroje } [\mathrm{bit/s}],\]

Zdrojové kódování tvrau vlny: pulzní modulace (PCM), diferenční modulace (DPCM)
\begin{itemize}
	\item Pulzní kódová PCM - analogový signál se: vzorkuje, kvantuje, každý vzorek se zakóduje do bitů
	\item Diferenční pulzní kódová DPCM - nepřenáší se celý vzorek, ale jen rozdíl mezi aktuálním vzorkem a jeho predikovanou hodnotou z předchozího vzorku
	\item Delta DM - přenáší pouze rozdíl, jestli se signál zvýšil/snížil oproti předchozímu vzorku (pouze 2 úrovně)
	\item Adaptivní delta ADM - princip DM, ale s proměnným kvantizačním krokem podle strmosti signálu (lépe sleduje rychlé změny)
	\item Adaptivní diferenční pulzní kódová ADPCM - princip DPCM, ale s adaptivním kvantováním i adaptivní predikcí (stejná kvalita jako PCM, při nižším datovém toku)
\end{itemize}

\subsection{Kanálové kódování}
Provádíme za účelem zvýšení odolnosti signálu proti chybám v přenosovém kanálu:
\begin{itemize}
	\item přidáním redundance pro detekci nebo opravu chyb
\end{itemize}
Např.: RS, LDPC; doplněné prokládáním

\subsection{Prokládání}
Ochrana signálu proti skupinovým chybám
\begin{itemize}
	\item Vysílač (kodér)
	\begin{itemize}
		\item 1) zpráva zakódovaná běžným kodérem
		\item 2) zápis zakódované zprávy do amtice M \textbf{po řádcích}
		\item 3) čtení symbolů \textbf{po sloupcích}
		\item 4) takto přeskládané symboly se vyšlou anténou
	\end{itemize}
	\item Přijímač (dekodér)
	\begin{itemize}
		\item 5) zápis přijatých symbolů opět \textbf{do sloupců} (rekonstrukce matice M)
		\item 6) čtení symbolů přijaté zprávy z matice M \textbf{po řádcích} $\rightarrow$ [přeskládání opět do původní podoby]
		\item 7) oprava případných chyb jednoduchými dekodéry
	\end{itemize}
\end{itemize}
Parametry prokládání:\\
\textit{Hloubka prokládání:} čím větší, tím delší shluk chybb dokáže systém eliminovat\\
\textit{Rámec vnějšího kódu:} udává počet bitů, pokterých se budou opakovat chybné bity

\subsection{Digitální modulace}
Proces, při kterém se digitální signál přenáší změnou parametrů nosné vlny.\\
Takto změněná nosná je anténou vysílače odeslána k přijímači.\\
\textbf{Obálka:} pomyslná křivka spojující vrcholky amplitud vysílané nosné vlny
\begin{itemize}
	\item \underline{lineární modulace:} proměnná obálka (informace skryta ve změnách amplitudy - AM, ASK)
	\item \underline{nelineární modulace:} konstantní obálka (informace skryta ve změnách frekvence nebo fáze - FM, PM, FSK, PSK)
\end{itemize}
Změna parametrů:
\begin{itemize}
	\item \textbf{ASK}: mění se skokově \textbf{amplituda} nosné vlny ("Amplitude Shift Keying")
	\item \textbf{PSK}: mění se skokově \textbf{fáze} nosné vlny ("Phase Shift Keying")
	\item \textbf{FSK}: mění se skokově \textbf{frekvence} nosné vlny ("Frequency Shift Keying")
	\item \textbf{QAM}: kombinuje současnou změnu \textbf{amplitudy} i \textbf{fáze}. Dokáže přenášet nezávislé signály n1, n2 stejným kanálem ("Kvadraturní AM")\\
\end{itemize}
Vícestavová modulace: 
\begin{equation}
	{M = 2^n},
\end{equation}
(\textit{M}: počet stavů nosné; \textit{n}: počet bitů jednoho slova)\\

\underline{\textbf{Modulace QPSK:}} 4 stavová fázová modulace
\begin{itemize}
	\item Každý fázový stav představuje jeden dibit (00, 01, 10, 11)
	\item Realizace složením dvou fázově posunutých BPSK
	\item Modulátor vytváří \textit{sin} větve Q a \textit{cos} větve I $\rightarrow$ vlny jsou fázově posunuté o 90°: 45°, 135°, -135°, -45°
	\item teoreticky by mělo jít o nelineární modulaci, v praxi ale vzniká parazitní amplitudová modulace při průchochodu středem diagramu (fázový posun o 180°)
	\item Oproti BPSK přenáší QPSK 2 bity na 1 symbol při stejné šířce pásma
	\item Změna stavu obou bitů jednoho symbolu ve stejný okamžik
\end{itemize}
\begin{figure}[H]
	\centering
	\includegraphics[width=0.2\linewidth]{"Obrazky/MKO/2. otazka/qpsk"}
	\caption{Vektorový diagram modulace QPSK}
	\label{fig:qpsk}
\end{figure}
\newpage
\underline{\textbf{Modulace O-QPSK:}} 4 stavová fázová modulace
\begin{itemize}
	\item časové zpoždění mezi změnami obou bitů $\rightarrow$ není průchod nulou
\end{itemize}
\begin{figure}[H]
	\centering
	\includegraphics[width=0.2\linewidth]{"Obrazky/MKO/2. otazka/oqpsk"}
	\caption{Vektorový diagram modulace O-QPSK}
	\label{fig:oqpsk}
\end{figure}


\underline{\textbf{Modulace $\pi/4$-QPSK:}} 8 stavová fázová modulace
\begin{itemize}
	\item používá 2 QPSK diagramy pootočené o 45°
	\item nenastane průchod 0 (posuv o 180°)
\end{itemize}
\begin{figure}[H]
	\centering
	\includegraphics[width=0.2\linewidth]{"Obrazky/MKO/2. otazka/pi4qpsk"}
	\caption{Vektorový diagram modulace $\pi/4$-QPSK}
	\label{fig:pi4qpsk}
\end{figure}


\underline{\textbf{Modulace MSK:}} 2 stavová frekvenční modulace ("Minimum shift keying")
\begin{itemize}
	\item změna fáze NENÍ skoková jako u předešlých modulací
	\item oba stavy nejsou od sebe fázově vzdáleny o 180° jako u BPSK, ale při přechodu se plynule mění o 90°
	\item "minimum": minimální frekvenční rozestup, při kterém jsou oba stavy ortogonální(pravoúhlý), aby se přenos dal rozlišit
\end{itemize}

\chapter{Třetí otázka}
\section{Systémy s mnohonásobným přístupem}
Princip: více nezávislých uživatelů odesílá k jednomu společnému médiu (např.: mobilní systém s mnoha uživateli)
\begin{itemize}
	\item FDMA (Freq Division Multiple Access): rozdělení kmitočtového pásma na jednotlivé kanály $\rightarrow$ každý uživatel dostane přidělenou svou vlastní frekvenci
	\item TDMA (Time Division Multiple Access): všichni uživatelé vysílají na stejné frekvenci, ale střídají se v čase $\rightarrow$ uživatelovi se přidělí krátký časový úsek, rychle se odešle zpráva a uživatel poté čeká na další kolo přidělování
	\item CDMA (Code Division Multiple Access): všichni uživatelé vysílají na stejné frekvenci i ve stejném čase; data jsou šifrována unikátními klíči 
	\item ALOHA: data se odešlou kdykoliv; pokud se 2 signály "srazí", vysílači přijme potvrzení o nepřijetí a data se odešlou znovu s náhodným časovým zpožděním znovu (přijatý kód musí mít správný formát zápisu, který se při srážce a spojení dvou signálů poškodí)
\end{itemize}

\section{Metody multiplexování}
Princip: z jednoho bodu se odesílá více různých signálů do jednoho radiového kanálu (např.: rádio s mnoha stanicemi) - odvozeno od mnohonásobného přístupu.
\begin{itemize}
	\item FDM (kmitočtový multiplex): rozdělení jednoho kanálu na několik frekvenčních pásem; signály se odesílají současně, ale každý na jiné frekvenci
	\begin{description}
		\item[$\oplus$] není nutná časová synchronizace, vhodné pro spojité signály
		\item[$\circleddash$] horší využití pásma kvůli ochranným intervalům, vyžaduje filtry
	\end{description}
	\item TDM (časový multiplex): všichni uživatelé používají stejný kanál, ale ne ve stejnou dobu; každý má přiřazen vlastní časový interval
	\begin{description}
		\item[$\oplus$] lepší využití kanálu než FDM (každý uživatel dostane celý kanál)
		\item[$\circleddash$] přesná synchronizace (může vznikat zpoždění)
	\end{description} 
	\item CDM (kódový multiplex): všichni uživatelé vysílají ve stejném čase i pásmu, ale každý pomocí jiného kódu
	\begin{description}	
		\item[$\oplus$] více uživatel sdílí stejné pásmo současně
		\item[$\circleddash$] složité zpracování, při vyšším počtu uživatel roste vzájemné rušení
	\end{description}
\end{itemize}
\begin{figure}[H]
	\centering
	\includegraphics[width=0.45\linewidth]{"Obrazky/MKO/3. otazka/multiplexovani"}
	\caption{Multiplexování ve frekvenční, časové a kódové doméně}
	\label{fig:multiplexovani}
\end{figure}

\section{Proces rozprostírání signálu a výhody/nevýhody}
Jev typický pro CDMA. Užitečný signál se násobí rozprostíracím šumovým kódem (např.: LSFR kódem - paralela s DE1) $\rightarrow$ umělé zvětšení šířky pásma a tím rozložení výkonu do širšího spektra. Na přijímači se stejný kód použije pro zpětné zúžení.\\
Důvody použití:
\begin{itemize}
	\item Potlačení úzkopásmového rušení
	\item Odolnost proti odposlechu a rušení
	\item Více uživatelů v jednom spektru
\end{itemize}

\chapter{Čtvrtá otázka}
\section{Způsoby přenosu}
\subsection{Simplex - jednosměrný přenos}
Jeden konec vysílač, druhý přijímač.\\
Bez možnosti otočení směru.\\ 
Např.: TV, FM radio
\subsection{Half-Duplex - střídavě obousměrný přenos}
Komunikace v bobou směrech, ale nikdy ne současně.\\
Např.: vysílačky push-to-talk
\subsection{Duplex - plně obousměrný přenos}
Komunikace v obou směrech současně.\\
Např.: telefonní hovor, kabelový internet
\section{Typy spojování 2+ účastníků sítě}
\subsection{Komutované}
Vytvoření spojení mezi účastníky $\rightarrow$ vyhrazené přenosové pásmo zaniká až při ukončení komunikace.
\begin{itemize}
	\item 1 přenosový kanál mezi dvěma účastníky
	\item vyžaduje požadovanou kvalitu spojení
	\item nižší efektivita využití přenosové kapacity kanálu (stále obsažený kanál)
	\item Při přenosu signálu se využívá \textit{synchronní přenosový mód} (s konstantní přenosovou rychlostí signálu)
\end{itemize}
\subsection{Paketové}
Tok data rozdělen na \textit{pakety} $\rightarrow$ ty se přenášejí samostatně a síť pro ně nedrží jednu, trvale rezervovanou, cestu.\\
Přenos probíhá způsoby:
\begin{itemize}
	\item Služba bez spojení - vhodná pro přenos krátkých zpráv
	\item Služba se spojením - první (vyhledávací) paket vyznačí \textit{virtuální okruh} a rezervuje v každém přepojovacím uzlu dostatečnou paměť pro celou zprávu. Pro velké objemy dat.
\end{itemize}
	\begin{description}	
	\item[$\oplus$] efektivní využití přenosové kapacity sítě
	\item[$\circleddash$] problémy vznikající při vzájemné komunikaci v reálném čase
\end{description}
\section{Diverzitní příjem}
Metoda využití více kanálů (cest) pro přenos jednoho signálu v případě výpadku jednoho kanálu (snížení chybovosti)\\
Příčiny výpadku signálu:
\begin{itemize}
	\item Únik (Kolísání signálu v místě příjmu) způsobený změnami úrovní terénu, překážkami mezi (pohybující se) mobilní stanicí  a (statickou) základnovou stanicí
	\item mnohacestným šířením mezi vysílačem a přijímačem, komplikovaným navíc ještě pohybem TX (= vysílač) a RX (= přijímač) $\rightarrow$ vektorové sčítání signálů
\end{itemize}
\begin{figure}[H]
	\centering
	\includegraphics[width=0.5\linewidth]{"Obrazky/MKO/4. otazka/diverzitniPrijem"}
	\caption{Vícecestný příjem}
	\label{fig:diverzitniprijem}
\end{figure}
Vytváření více kanálů:
\begin{itemize}
	\item \textbf{Prostorová diverzita} - použití více antén vzájemně vzdálených (o několik vlnových délek, min. $\lambda$/2)
	\item \textbf{Úhlová diverzita} - směrové antény nastavené do různých směrů (příjem signálu odražených od více objektů)
	\item \textbf{Polarizační diverzita} - 2 antény s různou horizontální a vertikální polarizací
	\item \textbf{Frekvenční diverzita} - posílání stejných dat na různých frekvencích (odstup jednotlivých kmitočtů 2-4 \% kmitočtu nosné vlny)
	\item \textbf{Časová diverzita} - posílání stejných dat s časovým odstupem (rámec TDM má délku až jednotky sekund - dlouhá doba přenosu)
\end{itemize}
Rozlišení diverzitních systémů podle počtu antén vysílače a přijímače:
\begin{itemize}
	\item \textbf{SISO} (single IN, single OUT) - 1 vysílá; 1 přijímá
	\item \textbf{SIMO} (single IN, multiple OUT) - 1 vysílá; více antén přijímá
	\item \textbf{MISO} (multiple IN, single OUT) - více antén vysílá; 1 přijímá
	\item \textbf{SISO} (multiple IN, multiple OUT) - více antén vysílá; více antén přijímá
\end{itemize}

\chapter{Pátá otázka}
\section{Struktura obecné buňkové sítě a hlavní výhody}
Buňková síť rozděluje celé pokrytí na \textit{buňky}, ve kterých pracují základnové stanice.\\
Každá buňka obsluhuje omezenou oblast a sousední buňky mohou znovu používat stejné kmitočty, pokud je zajištěn dostatečný odstup, aby nedocházelo k rušení.
\subsection{Základní princip:}
Mobilní stanice komunikuje s nejbližší základnovou stanicí, která zprostředkuje přístup do vyšší vrstvy sítě.\\
Síť je obvykle tvořena:
\begin{itemize}
	\item mobilními stanicemi uživatelů,
	\item základnovými stanicemi,
	\item řídicí a přepojovací částí sítě,
	\item rozhraním do dalších sítí.
\end{itemize}
\subsection{Plošná struktura}
\begin{itemize}
	\item \textbf{Spojení bod - bod} (PP point to point) - pro spojení ústředny s \textbf{jedním} místem, kde je vyšší koncentrace účastníků. Realizace mikrovlnným spojem se směrovými anténami
	\item \textbf{Spojení bod - více bodů} (PM point to multipoint) - pro spojení ústředny s \textbf{několika} místy vyšší koncentrace účastníků, na které mohou navazovat bezdrátová spojení jednotlivých účastníků
	\begin{figure}[H]
		\centering
		\includegraphics[width=0.5\linewidth]{"Obrazky/MKO/5. otazka/bunkovaStruktura"}
		\caption{Buňková struktura sítě s propojením bod-více bodů}
		\label{fig:bunkovastruktura}
	\end{figure}
	\item \textbf{Buňková struktura} - území rozděleno na velké množství menších oblastí (buněk). Uvnitř buňky je umístěna základnová radiová stanice. Buňky mají různou velikost i tvar:
	\begin{itemize}	
		\item \textbf{femtobuňka} (poloměr desítky \textit{m}) - pokrývá nepokryté části území mezi menšími piko nebo mikrobuňkami
		\item \textbf{pikobuňka} (poloměr $<$50 \textit{m}) - pro místa s vysokou koncentrací uživatelů (nádraží, obchodní domy)
		\item \textbf{mikrobuňka} (poloměr $<$1 \textit{km}) - oblasti s větším provozem (centra měst)
		\item \textbf{makrobuňka} (poloměr desítky \textit{km}) - oblasti s malou hustotou provozu (venkov)
	\end{itemize}
\end{itemize}
	\begin{figure}[H]
		\centering
		\includegraphics[width=0.45\linewidth]{"Obrazky/MKO/5. otazka/bunkovaStruktura2"}
		\caption{Ukázka systému buněk}
		\label{fig:bunkovastruktura2}
	\end{figure}
	Buňky očíslovány 1 až 7 v rozsahu 1-100 až 601-700; každá buňka používá jiný kmitočtový kanál a stejná čísla se tedy mohou opakovat, pokud jsou v dostatečně vzdálené oblasti\\
	\underline{\textit{Stejné číslo}} = stejná frekvenční sada v rámci opakující se frekvenční struktury\\
	\underline{\textit{Jiná čísla}} = jiné kanálové skupiny, aby sousední buňky nerušily jedna druhou\\
	Stejné frekvence se použijí vícekrát v síti, ale jen v buňkách, které jsou od sebe dost daleko $\rightarrow$ výrazné zvýšení kapacity systému.\\
	\textbf{Zvýšení kapacity sítě:} Jedna buňka může mít 3 nebo 6 sektorů $\rightarrow$ zvýšení kapacity a zabránění rušení pokrýváním buňky po částech pomocí sektorových antén. 
	
\subsection{Hlavní výhody buňkové sítě}
\begin{itemize}
	\item \textbf{Vysoká kapacita} díky opakovanému použití frekvencí v různých buňkách.
	\item \textbf{Úspora kmitočtového spektra}, protože není nutné mít pro každého účastníka vlastní frekvenci.
	\item \textbf{Dobré pokrytí} i rozsáhlého území pomocí sítě menších buněk.
	\item \textbf{Podpora mobility} uživatele při pohybu mezi buňkami.
	\item \textbf{Možnost zvyšovat kapacitu} zmenšováním buněk, například ve městech.
\end{itemize}

\section{nevyžaduje, ale pro úplnost: Nevýhody a omezení}
\begin{itemize}
	\item Nutnost řízení rušení mezi sousedními buňkami.
	\item Složitější správa sítě.
	\item Potřeba předávání spojení při pohybu uživatele mezi buňkami.
\end{itemize}

\section{Handover}
Proces, při kterém se mobilní spojení přepojí na jinou mobilní stanici (bez přerušení komunikace), když se uživatel pohybuje mezi buňkami. Kdyby nedošlo k předání spojení, mobilní stanice by při opuštění buňky ztratila spojení.\\
\textbf{Průběh:}
\begin{enumerate}
	\item Systém sleduje kvalitu signálu od sousedních buněk.
	\item Vyhodnotí, že současná buňka už není vhodná pro další obsluhu.
	\item Vybere se cílová základnová stanice.
	\item Přenos se předá nové stanici.
	\item Původní spojení se ukončí.
\end{enumerate}
\subsection{Typy handoveru}
\begin{itemize}
	\item \textbf{Tvrdý handover} odpojení od 1. buňky $\rightarrow$ přerušení $\rightarrow$ připojení k 2. buňce
	\begin{figure}[H]
		\centering
		\includegraphics[width=0.4\linewidth]{"Obrazky/MKO/5. otazka/handoverTvrdy"}
		\caption{Tvrdý handover}
		\label{fig:handovertvrdy}
	\end{figure}
	\item \textbf{Měkký handover} mobilní telefon komunikuje současně s více buňkami (základnovými stanicemi), než definitivně opustí signál té původní
	\begin{figure}[H]
		\centering
		\includegraphics[width=0.3\linewidth]{"Obrazky/MKO/5. otazka/handoverMekky"}
		\caption{Měkký handover}
		\label{fig:handovermekky}
	\end{figure}
	\item \textbf{Bezešvý handover} krátkodobé překrytí spojení při přepojení mezi kanály
	\begin{figure}[H]
		\centering
		\includegraphics[width=0.4\linewidth]{"Obrazky/MKO/5. otazka/handoverBezesvy"}
		\caption{Bezešvý handover}
		\label{fig:handoverbezesvy}
	\end{figure}
\end{itemize}
\section{Rozdělení handoveru podle způsobu měření kvality spojení}
$\rightarrow$ rozhodují o handoveru a řídí jej
\begin{itemize}
	\item \textbf{Sítí řízený handover} měření kvality kanálu provádí základnová stanice (na základě výsledků rozhoduje o přepnutí)
	\begin{description}	
		\item[$\oplus$] na mobilní stanici (UE - user equipment) nejsou kladeny žádné požadavky
	\end{description}
	\item \textbf{Handover řízený mobilní stanicí} měření kvality provádí mobilní stanice i základnová stanice (rozhodnutí o přepnutí provede mobilní stanice, předá jej do systému a ten provede přepnutí) 
	\item \textbf{Sítí řízený handover s asistencí mobilní stanice} mobilní stanice měří velikost signálu sousedních základnových stanic a výsledky měření předává servisní základnové stanici (ke které je právě připojena).\\
	Současně mobilní i základnová stanice provádí měření probíhajícího spojení. Na základě naměřených údajů provádí systém rozhodnutí o přepnutí spojení a uskuteční je (např. GSM).
\end{itemize}
\chapter{Šestá otázka}
\section{Základní architektura a vlastnosti přenosových systémů GSM, GPRS, EDGE, UMTS}
\subsection{GSM 2G}
Bezdrátový systém mobilních sítí GSM navržen tak, aby nebyl uzavřený, ale aby umožňoval přístup i do jiných sítí (pevná telefonní linka, mobilní sítě, např. LTE).\\
\underline{Přenos:}  spojení trvale rezervované, i když uživatel zrovna nic neposílá.\\
\underline{Architektura:}
\begin{itemize}
	\item \textbf{BSS (base station subsystem)} subsystém základnových stanic BTS; komunikace s UE přes rádiové rozhraní BSC.
	\item \textbf{NSS (network \& switching subsystem)} síťový a přepojovací subsystém - směrování hovorů do jiných sítí
	\item \textbf{OSS (operation support subsystem)} operační provozní subsystém - monitoring, údržba, správa parametrů sítě
\end{itemize}
\begin{figure}[H]
	\centering
	\includegraphics[width=0.4\linewidth]{"Obrazky/MKO/6. otazka/architekturaGSM"}
	\caption{Architektura GSM}
	\label{fig:architekturagsm}
\end{figure}
\subsection{GPRS}
Pomocí systému GPRS je rozšířen GSM a umožňuje \textbf{přenos datových paketů} přes rádiové rozhraní s teoretickou přenosovou rychlostí max 171,2 kbit/s.\\
\underline{Paketový přenos:} rádiozdroje se přidělují jen tehdy, když je co posílat (= lepší využití kapacity sítě)
Doplnění stávající GSM sítě o bloky:\\
\begin{itemize}
	\item \textbf{PCU (Packet controller unit)} zajišťuje řízení paketového provozu na rádiovém rozhraní.
	\item \textbf{SGSN (Serving GPRS support node)} provádí šifrování dat a kompresi; komunikace s rádiovou částí sítě GPRS a druhým datovým uzlem GGSN
	\item \textbf{GGSN (Gateway GPRS support node)} směrovač, jehož úkolem je komunikace s externími paketovými datovými sítěmi (např. sítí Internet)
\end{itemize}
\begin{figure}[H]
	\centering
	\includegraphics[width=0.4\linewidth]{"Obrazky/MKO/6. otazka/architekturaGPRS"}
	\caption{Architektura GPRS}
	\label{fig:architekturagprs}
\end{figure}

\subsection{EDGE}
Další evoluční stupeň GSM/GPRS\\
Pomocí efektivní digitální modulace 8-PSK (namísto GMSK) zvyšuje přenosovou rychlost systému GSM.\\
3x vyšší přenosová rychlost než GPRS (GPRS: 270 kbit/s, EDGE: 810 kbit/s)\\
Díky modulaci 8-PSK je náchylnější na přenosové podmínky.
\begin{figure}[H]
	\centering
	\includegraphics[width=0.2\linewidth]{"Obrazky/MKO/6. otazka/8psk"}
	\caption{Vektorový diagram modulace 8PSK}
	\label{fig:8psk}
\end{figure}

\subsection{UMTS 3G}
Oproti GSM/GPRS/EDGE přenáší výrazně vyšší kapacitu, lepší datové služby a širší možnost mobilního přístupu k síti.\\
Architektura:
\begin{itemize}
	\item \textbf{RNS (Rádiová přístupová síť)} využívá pozemní rádiové rozhraní UTRAN (universal terrestrial radio access network) - UE, NodeB, RNC (= přenosová a přepojovací funkce)
	\item \textbf{CN (páteřní síť)} řídí provoz a spojení v systému
\end{itemize}
Jednotlivé prvky architektury:
\begin{itemize}
	\item \textbf{UE (user equipment)} - mobilní stanice MS systému UMTS může pracovat v módech:
	\begin{itemize}	
		\item CS (circuit switch mode pro hlasové hovory)
		\item PS (packet switch mode pro datové přenosy)
		\item PS/CS (circuit switch mode pro připojení současně k PS i CS)
	\end{itemize}
	\item \textbf{NodeB (základnová stanice u 3G)}
	\begin{itemize}	
		\item převod signálu přes antény z transportních kanálů do fyzických kanálů
		\item řízení výkonu mobilních stanic (důležité pro handover)
		\item kmitočtová a časová synchronizace (chipová, bitvá, slotová, rámcová)
	\end{itemize}
	\item \textbf{RNC (kontrolér)}
	\begin{itemize}	
		\item zajišťuje řízení UTRAN
		\item vysílá systémové informace o buňkách
		\item řídí handover 
	\end{itemize}
	\item \textbf{CS (circuit switched) Domain} - komutovaný provoz s přepínáním okruhů
	\begin{itemize}	
		\item pro hlasovou komunikaci
	\end{itemize}
	\item \textbf{PS (packet switched) Domain} - paketový provoz s přepínáním paketů
	\begin{itemize}	
		\item pro datový přenos (např. se sítí Internet)
	\end{itemize}
	\item \textbf{BC (broadcast) Domain}
	\begin{itemize}	
		\item centrum pro koordinaci vysílání v jednotlivých buňkách CNC (cell broadcast center)
	\end{itemize}
\end{itemize}

\begin{figure}[H]
	\centering
	\includegraphics[width=0.5\linewidth]{"Obrazky/MKO/6. otazka/architekturaUMTS"}
	\caption{Architektura UMTS}
	\label{fig:architekturaumts}
\end{figure}
Rozhraní mezi prvky systému UMTS:
\begin{itemize}
	\item \textbf{Uu} rozhraní mezi mobilním telefonem (UE) a vysílačem (NodeB)
	\item \textbf{Iub} rozhraní mezi vysílačem (NodeB) a kontrolérem (RNC)
	\item \textbf{Iur} rozhraní spojující více kontrolérů (RNC)
	\item \textbf{Iu} rozhraní mezi kontrolérem (RNC) a páteřní sítí (CN)
\end{itemize}
\begin{figure}[H]
	\centering
	\includegraphics[width=0.5\linewidth]{"Obrazky/MKO/6. otazka/architekturaUMTS-rozhrani"}
	\caption{Rozhraní sítě UMTS}
	\label{fig:architekturaumts-rozhrani}
\end{figure}

\chapter{Sedmá otázka}
\section{Systém LTE 4G}
Spadá mezi standardy UMTS.\\
Vysokorychlostní bezdrátový přenos dat.\\
Cílem LTE: zvýšení přenosové rychlosti, zkrácení odezvy, snížení nákladů\\
Šířka pásma: 1,4 až 20 MHz (flexibilní pro práci v různých RF pásmech)\\
Podpora MIMO: využití více vysílacích i přijímacích antén pro současný přenos více nezávislých datových toků na stejné frekvenci.\\
Podpora obou duplexů:
\begin{itemize}
	\item \textbf{Frequency division duplex (FDD)}: downlink i uplink probíhají současně, každý ale odděleně na jiné frekvenci
	\item \textbf{Time division duplex (TDD)} downlink i uplink probíhají na stejné frekvenci, ale velmi rychle se střídají v čase
\end{itemize}
\subsection{Modulace:}
\begin{itemize}
	\item \textbf{Vnitřní (mapování datových bitů na symboly)}
	\begin{itemize}	
		\item QPSK: nejodolnější modulace pro špatný signál; přenese 2bity/symbol
		\item 16-QAM: pro středně dobrý signál; přenese 4 bity/symbol
		\item 64-QAM: pro čistý signál; přenese 6 bitů/symbol (nejrychlejší v základním LTE)
	\end{itemize}
	\item \textbf{Vnější (symboly rozděleny do navzájem kolmých subnosných frekvencí)} 
	\begin{itemize}	
		\item OFDM (rozdělení širokého 20 MHz spektra na několik subnosných)
	\end{itemize}
\end{itemize}

\subsection{Architektura:}
\begin{itemize}
	\item \textbf{UE}: koncový uživatel, bezdrátové spojení se sítí rozhraním Uu
	\begin{itemize}	
		\item SIM karta (slouží k identifikaci)
		\item Mobilní zařízení (slouží k připojení do sítě)
		\item Terminál TE (hardware: displej, reproduktor, mikrofon,...)
	\end{itemize}
	\item \textbf{E-UTRAN}: přístupová rádiová síť, skládá se ze základnových stanic eNodeB
	\begin{itemize}	
		\item vytváří rádiové buňky
		\item nepotřebuje kontrolér RNC, protože jsou stanice eNodeB propojeny mezi sebou přímo $\rightarrow$ plynulé přepojení při handoveru
	\end{itemize}
	\item \textbf{EPC}: jádro sítě
	\begin{itemize}	
		\item plně paketovaná síť operátora
		\item řídí síť, ověřuje uživatele, přiděluje IP adresy
	\end{itemize}
\end{itemize}
\begin{figure}[H]
	\centering
	\includegraphics[width=0.5\linewidth]{"Obrazky/MKO/7. otazka/architekturaLTE"}
	\caption{Architektura LTE}
	\label{fig:architekturalte}
\end{figure}
Jednotlivé prvky architektury:
\begin{itemize}
	\item \textbf{eNodeB} - základnová stanice (anténa), řídící jednotka rádiové sítě
	\begin{itemize}	
		\item modulace
		\item rozdělování pásem uživatelům
		\item řídí handover 
	\end{itemize}
	\item \textbf{MME} - hlavní řídící prvek sítě LTE
	\begin{itemize}	
		\item obsluha několika eNodeB
		\item kontrola přístupu do sítě
		\item ochrana proti odposlechu
	\end{itemize}
	\item \textbf{SAE-GW (dvojice servisní brána S-GW a brána paketové a datové sítě P-GW)} 
	\begin{itemize}	
		\item udržuje datový tok, připojuje k Internetu, přiděluje IP adresy
	\end{itemize}
\end{itemize}
Rozhraní mezi prvky systému UMTS:
\begin{itemize}
\item \textbf{Uu} - bezdrátové rádiové rozhraní mezi UE a eNodeB
\item \textbf{X2} - kabelové rozhraní mezi základnovými stanicemi eNodeB
\end{itemize}

\subsection{LTE-Advanced}
Standard pro systémy 4G.\\
\begin{itemize}
	\item Latence: 5 ms
	\item 3x vyšší spektrální účinnost (= datový tok/šířka pásma) oproti LTE (downlink: 30bit/s/Hz; uplink: 15bit/s/Hz)
	\item více antén (více uživatelů přenáší současně)
	\begin{itemize}	
		\item 8x8 MIMO pro downlink
		\item 4x4 MIMO pro uplink
	\end{itemize}
	\item sdružení sousedních i nesousedních nosných pásem do jednoho výsledného $\rightarrow$ velmi rychlá přenosová rychlost
	\begin{itemize}	
		\item sloučením dvou 20 MHz pásem vznikne pásmo 40 MHz (v kombinaci s 8x8 MIMO umožňuje přenosovou rychlost v downlinku 1,2 Gbit/s)
		\item maximálně lze sloučit 5 pásem (max výsledná šířka pásma: 100 MHz, přenosová rychlost downlink: 3 Gbit/s, přenosová rychlost uplink: 1,5 Gbit/s)
	\end{itemize}
	\item Využití Reléových stanic RN
	\begin{itemize}	
		\item zvýšení kapacity sítě i zlepšení kvality spojení mezi eNodeB a UE
		\item zvětšení pokrytí území pomocí \textbf{femtobuněk}
		\begin{itemize}	
			\item pro lokální pokrytí v místech se špatnou kvalitou signálu/
			\item v hustě osídlených lokalitách pro vysoké přenosové rychlosti
		\end{itemize}
		\begin{figure}[H]
			\centering
			\includegraphics[width=0.6\linewidth]{"Obrazky/MKO/7. otazka/femtobunky"}
			\caption{Aplikace femtobuněk (vlevo: hustě osídlené území, vpravo: v místě se špatným signálem-mimo dosah)}
			\label{fig:femtobunky}
		\end{figure}
		
	\end{itemize}
\end{itemize}

\chapter{Osmá otázka}
\section{Systém LTE 4G}
\subsection{E-UTRA Downlink}
Použití systému OFDMA (více uživatelů sdílí stejné spektrum, každý dostane svou podmnožinu subnosných) pro vyšší spektrální účinnost ve srovnání s WCDMA (rozkládání uživatelských bitů do širšího frekvenčního pásma; u 3G)
\begin{figure}[H]
	\centering
	\includegraphics[width=0.7\linewidth]{"Obrazky/MKO/8. otazka/downlinkOFDM"}
	\caption{Znázornění OFDM}
	\label{fig:downlinkofdm}
\end{figure}
\begin{figure}[H]
	\centering
	\includegraphics[width=0.7\linewidth]{"Obrazky/MKO/8. otazka/downlinkOFDMA"}
	\caption{Znázornění OFDMA} (rozestupy subnosných 15 kHz)
	\label{fig:downlinkofdma}
\end{figure}
\begin{itemize}
	\item \textbf{Princip:} Data jsou přenášena paralelně na velkém množství úzkopásmových ortogonálních subnosných.
	\item \textbf{Modulace:} Používá se QPSK, 16-QAM nebo 64-QAM.
	\item \textbf{Výhody:} Vysoká spektrální účinnost, dobrá odolnost vůči mezisymbolovým interferencím (ISI) díky ochrannému intervalu (\textit{CP Cyclic Prefix}), snadná implementace MIMO.
	\item \textbf{Nevýhody:} Značné kolísání obálky signálu, což vede k vysokému PAPR (Poměr maximálního s středního výkonu signálu). To ale nevadí základnové stanici (eNodeB), která má dostatek napájecího výkonu.
	\item Zdrojové prostředky (RB) jsou přidělovány uživatelům na základě jejich požadavků, aktuálním stavu sítě a časové kmitočtové mapy úniků
	\item Nejmenší kmitočtovou časovou jednotkou je zdrojový blok RB (resource block)
	\item Nejmenší přenosovou jednotkou je SB (scheduling block), tvořený dvěma PRB, které na sebe bezprostředně časově navazují (12 subnosných $\times$ 0,5 ms)
	\item Jednomu uživateli může být přiděleno i více zdrojových bloků, které mohou být rozprostřeny v dostupném kmitočtovém pásmu
	\item V 1. a 5. OFDM symbolu v každém slotu se přenáeji referenční signály (v DL specifikují buňku, UE a službu MBSFN (= způsob vysílání, kdy více vysílačů vysílá stejný signál na stejné frekvenci, pro lepší pokrytí))
	\begin{itemize}	
		\item \textbf{Referenční signál} - pro odhad charakteristiky kanálu (UE podle nich měří sílu přenosu a rušení); využití v DL i UL
		\item \textbf{Synchronizační signál} - identifikační buňky (synchronizace UE se sítí, aby mohla začít komunikace)
		\begin{figure}[H]
			\centering
			\includegraphics[width=0.25\linewidth]{"Obrazky/MKO/8. otazka/refSymbolyDownlink"}
			\caption{Umístění referenčních symbolů v LTE (downlink)}
			\label{fig:refsymbolydownlink}
		\end{figure}
	\end{itemize}
\end{itemize}
	\newpage
	Proces v downlinku:
\begin{itemize}
	\item (1) ochrana proti chybám + prokládání dat
	\item (2) tvorba datových symbolů modulací (QPSK, 16-QAM, 64-QAM)
	\item (3) převod do paralelního tvaru a rozdělení na jednotlivé subnosné
	\item (4) vytvoření OFDM signálu, tedy převod z kmitočtové do časové oblasti
	\item (5) sloučení signálu + vložení ochranného intervalu (cyklický prefix)
	\item (6) převod na analogový signál + transpozice na vysílací nosnou před odesláním anténou
	
	\begin{figure}[H]
		\centering
		\includegraphics[width=0.7\linewidth]{"Obrazky/MKO/8. otazka/downlinkSchema"}
		\caption{Blokové schéma pro E-UTRA (Downlink)}
		\label{fig:downlinkschema}
	\end{figure}
\end{itemize}		

\subsection{E-UTRA Uplink}
Použití systému SC-FDMA (\textit{Single Carrier - Frequency Division Multiple Access}), kvůli nízkému poměru špičkového a průměrného výkonu (UE nepracuje s velkými výkyvy výkonu při uploadu)
\begin{itemize}
	\item \textbf{Princip:} Před samotným OFDM mapováním je signál transformován pomocí diskrétní Fourierovy transformace (DFT), čímž se rozprostře přes přidělené subnosné. Z toho důvodu je signál vnímán spíše jako signál o jedné nosné (Single Carrier).
	\item \textbf{Výhody:} Výrazně nižší PAPR oproti OFDMA $\rightarrow$ efektivnější a levnější provedení koncového výkonového zesilovače v mobilním telefonu $\rightarrow$ prodlužuje výdrž baterie.\\
\end{itemize}
	
Proces v uplinku:
\begin{itemize}
	\item (1) ochrana proti chybám + tvorba symbolů
	\item (2) zpracování symbolů (omezení výkonových špiček)
	\item (3) určení, na jakých frekvencích bude UE vysílat do eNodeB
	\item (4) a (5) převedení do časové oblasti
	\item (6) převod na analog. + finální úpravy před odesláním anténou

	\begin{figure}[H]
		\centering
		\includegraphics[width=0.7\linewidth]{"Obrazky/MKO/8. otazka/uplinkSchema"}
		\caption{Blokové schéma pro E-UTRA (Uplink)}
		\label{fig:uplinkschema}
	\end{figure}
\end{itemize}
\subsection{Systém LTE-CAT}
Rozšířené LTE určené pro \textbf{Internet věcí (IoT) a M2M (Machine-to-Machine)} komunikaci.\\
	\begin{description}	
	\item[$\oplus$] výborný dosah signálu
	\item[$\oplus$] úsporné režimy (fungování s malou baterií i několik let)
	\item[$\oplus$] plynulý přenos dat při handoveru
\end{description}
\newpage

\begin{itemize}	
	\item \textbf{LTE Cat-M1}
	\begin{itemize}	
		\item \textbf{Parametry:} Šířka pásma je omezena na \textbf{1,4 MHz}
		\item \textbf{Rychlost:} Okamžitá maximální rychlost je okolo 1 Mbps v DL i UL.
		\item \textbf{Energetická účinnost:} Zařízení mají navržené spánkové mechanismy, čímž baterie vydrží i několik let.
		\item \textbf{Pokrytí:} Vyšší pokrytí oproti klasickému LTE (zisk cca o 13 dB).
		\item \textbf{Příklady:} chytré hodinky, GPS, zabezpečovací systémy (malý přenos dat)
	\end{itemize}
	\item \textbf{LTE Cat-M2 (nástupce M1)}
	\begin{itemize}	
		\item \textbf{Parametry:} Šířka pásma \textbf{5 MHz}
		\item \textbf{Rychlost:} Okamžitá maximální rychlost je okolo 4$\div$7 Mbps.
		\item \textbf{Příklady:} kamerové systémy s nižším rozlišením, chytré billboardy (větší přenos dat)
	\end{itemize}
\end{itemize}
\subsection{Systém NB-IoT}
Síť, kde je potřeba přenos malého množství dat při nízké přenosové rychlosti (desítky$\div$stovky kbps) s vysokým zpožděním.\\
Úzkopásmová bezdrátová technologie určená pro připojení velkého množství jednoduchých zařízení.
\begin{description}	
	\item[$\oplus$] Nízká cena
	\item[$\oplus$] Realizace pouze softwarovou úpravou LTE vysílačů
	\item[$\oplus$] Obsluha desítek tisíc zařízení z 1 buňky
	\item[$\oplus$] Úspora baterie (výdrž až 10 let)
\end{description}
\begin{itemize}	
	\item \textbf{NB-IoT scénáře:}
	\begin{itemize}	
		\item \textbf{Stand-alone:} NB-IoT běží samostatně na vlastní vyhrazené frekvenci (například na staré vyřazené 2G)
		\begin{figure}[H]
			\centering
			\includegraphics[width=0.3\linewidth]{"Obrazky/MKO/8. otazka/NB-IoTstandAlone"}
			\caption{NB-IoT Stand-alone}
			\label{fig:nb-iotstandalone}
		\end{figure}
		\item \textbf{Guard Band:} využití volné mezery v ochranném pásmu (NB-IoT je úzkopásmové a vejde se do hluchého místa)
		\begin{figure}[H]
			\centering
			\includegraphics[width=0.3\linewidth]{"Obrazky/MKO/8. otazka/NB-IoTguardBand"}
			\caption{NB-IoT Guard Band}
			\label{fig:nb-iotguardband}
		\end{figure}
		\item \textbf{In-Band:} využití jednoho z bloků uprostřed aktuálně používané LTE frekvence (NB-IoT se vysílá společně s LTE daty)
		\begin{figure}[H]
			\centering
			\includegraphics[width=0.3\linewidth]{"Obrazky/MKO/8. otazka/NB-IoTinBand"}
			\caption{NB-IoT In Band}
			\label{fig:nb-iotinband}
		\end{figure}
	\end{itemize}
\end{itemize}
\newpage
\chapter{Devátá otázka}
\section{Systémy WLAN}
Bezdrátová lokální síť určená pro připojení zařízení na kratší vzdálenost
\subsection{Standard IEEE 802.11 - WiFi}
Mezinárodní sada norem, definující způsob bezdrátové komunikace.\\
Původní standard 802.11 navržen jako bezdrátová alternativa k Ethernetu.\\
Varianty se liší: rychlostí přenosu, efektivitou, podporou víceuživatelského provozu
\begin{figure}[H]
	\centering
	\includegraphics[width=0.7\linewidth]{"Obrazky/MKO/9. otazka/IEEE"}
	\caption{IEEE normy}
	\label{fig:ieee}
\end{figure}

\subsection{WiFi - Fyzické vrstvy}
Systém, který zajišťuje převod bitů na radiový signál a zpět (volba pásma, kanálu, modulace, rychlosti přenosu)\\
Metody přenosu dat na fyzické vrstvě WiFi:
\begin{itemize}	
	\item \textbf{FH-SS (frequency hopping spread spectrum)}
	\begin{itemize}	
		\item rozšíření spektra frekvenčním "skákáním" podle pseudonáhodné sekvence $\rightarrow$ vysílač nevysílá na jedné frekvenci, ale nesutále velmi rychle přeskakuje mezi frekvencemi v rozsahu RF pásma (šířka pásma: 1 MHz)
		\item odolné proti rušení
		\item nízká přenosová rychlost (použití také v Bluetooth)
		\item Vyřazena z nových standardů
	\end{itemize}
	\item \textbf{DS-SS (direct sequence spread spectrum)}
	\begin{itemize}	
		\item Rozprostření spektra pseudonáhodnou sekvencí za použití logické funkce XOR na šířku pásma 20 MHz (signál vypadá jako šum - je odolný proti úzkopásmovému rušení a bez znalosti dešifrovacího PN kódu je nemožné signál odposlouchávat)
		\item Přijímač zachytí tento "šum" a vynásobí jej PN kódem (synchronizovaným s vysílačem), tím se signál opět zúží do původního čitelného stavu
		\item Rychlost přenosu při 802.11: 2 Mbit/s; při 802.11b: 11Mbit/s s DS-SS
	\end{itemize}
	\item \textbf{OFDM (Orthogonal frequency division multiplexing)}
	\begin{itemize}	
		\item spektrum přenosového pásma rozděleno na velké množství subnosných, vzájemně se nerušících, paralelních signálů
		\item použití u IEEE 802.11a (5 GHz),  IEEE 802.11g (54 Mbit/s)
	\end{itemize}
\end{itemize}
\subsection{WiFi - Předcházení kolizím}
Přístup pomocí paketů \textbf{RTS} (request to send) a \textbf{CTS} (clear to send).\\
Pro přenos každého paketu se používají 4 rámce:
\begin{itemize}	
	\item 1. RTS - požadavek na vysílání (PC $\rightarrow$ Router)
	\item 2. CTS - pokud kanál volný, připravený k vysílání (Router $\rightarrow$ PC)
	\item 3. Odeslání dat s krátkým delayem (PC $\rightarrow$ Router)
	\item 4. Čekání na potvrzení příjmu (Router $\rightarrow$ PC); potvrzení nepřijde $\rightarrow$ odeslání znovu
	\begin{figure}[H]
		\centering
		\includegraphics[width=0.45\linewidth]{"Obrazky/MKO/9. otazka/kolize"}
		\caption{WiFi - Předcházení kolizím}
		\label{fig:kolize}
	\end{figure}
\end{itemize}
\section{Systémy WPAN}
WPAN (\textit{Wireless Personal Area Network}) je bezdrátová osobní síť pro komunikaci na velmi krátkou vzdálenost, typicky v řádu několika metrů až desítek metrů.\\
Použití pro propojení osobních zařízení (mobil, sluchátka, klávesnice, myš).
\begin{itemize}
	\item Využití FH-SS 
	\item RF pásmo: 2,4 GHz
\end{itemize}
\subsection{Bluetooth}
\begin{itemize}
	\item Práce v bezlicenčním pásmu 2,4 GHz (stejně jako WiFi)
	\item modulace GFSK (Gauss. Freq. Shift Keyking) - před klíčováním je signál "vyhlazen" Gaussovým filtrem pro potlačení rušení
	\item počet kanálů: 79; šířka kanálu 1 MHz
	\item \textbf{RSSI:} indikátor síly signálu [dBm], pro účely lokalizace $\rightarrow$ vyšší hodnota $\rightarrow$ silnější signál
	\item \textbf{BLE:} (režim Bluetooth Low Energy): šířka pásma: 2 MHz; počet kanálů: 40
	\item \textbf{Topologie ad hoc:} využití architektury založené na \textit{Master} a \textit{Slave} uspořádání malých siťových struktur: \textbf{pikonetů} (pikosítí)
	\begin{itemize}	
		\item Komunikace: \textbf{bod-bod} nebo \textbf{bod-více bodů}
		\item Komunikace probíhá přes \textit{Mastera} k zařízení \textit{Slave} (\textit{Slave} zařízení spolu nemohou komunikovat napřímo)
		\item \textit{Master} určuje: identifikaci účastníků a jejich vzájemnou synchronizaci
		\item \textit{Slave} odpovídá na výzvy (určuje obsah dat odeslaný směrem k \textit{Master})
		\item Příklad takové komunikace: chytré hodinky (Slave) uloží data ze senzoru do paměti a odešlou je do mobilu (Master) na vyžádání
		\item Funkce zanikají s ukončením spojení (při následující komunikaci může funkci řídící jednotky plnit jiný terminál)
	\end{itemize}
\begin{figure}[H]
	\centering
	\includegraphics[width=0.5\linewidth]{"Obrazky/MKO/9. otazka/WPANtopologie"}
	\caption{Topologie ad hoc (vlevo), Komunikace mezi více ad hoc (vpravo)}
	\label{fig:wpantopologie}
\end{figure}
	\item \textbf{Bluetooth paket:}
	\begin{itemize}	
		\item Přístupový kód: sycnhronizace a autorizace přístupu do sítě (unikátní pro každou \textbf{pikosíť})
		\item Záhlaví: řídící informace (typ paketu, adresa spojení se sítí \textbf{pikonet})
		\item Data: uživatelská data, která chceme přenést
		\begin{figure}[H]
			\centering
			\includegraphics[width=0.7\linewidth]{"Obrazky/MKO/9. otazka/WPANpakety"}
			\caption{Bluetooth paket}
			\label{fig:wpanpakety}
		\end{figure}
	\end{itemize}
\end{itemize}

\chapter{Desátá otázka}
\section{Systém LPWAN}
LPWAN (\textit{Low Power Wide Area Network}) označuje skupinu bezdrátových sítí navržených pro \textbf{malé datové objemy, dlouhý dosah a velmi nízkou spotřebu energie}.\\
\textbf{Vlastnosti LPWAN:}
\begin{itemize}
	\item nízká spotřeba energie,
	\item malý datový tok (krátké zprávy),
	\item velký dosah,
	\item velký počet koncových zařízení,
	\item nízké provozní náklady,
	\item vhodnost pro senzory a jednoduchá IoT zařízení.
\end{itemize}

\subsection{IoT}
IoT (\textit{Internet of Things}) je koncept propojení fyzických zařízení, senzorů a objektů do sítě, aby mohla sbírat, odesílat a zpracovávat data.\\
Typické IoT aplikace jsou měření teploty, vlhkosti, spotřeby energie, sledování polohy nebo průmyslová automatizace.\\
LPWAN technologie se pro IoT hodí právě proto, že umožňují levné, energeticky úsporné a dostatečně spolehlivé připojení mnoha zařízení.
\begin{itemize}
	\item \textbf{IoT}
	\begin{itemize}
		\item Platforma počítačů, chytrých zařízení a strojů schopných vzájemné komunikace a spolupráce.
		\item Vzájemné odesílání a příjem dat.
		\item IoT zařízení má vlastní cloud, software, senzory a možnost síťové konektivity.
		\item Princip fungování:
		\begin{itemize}
			\item senzory přijmou signál,
			\item proběhne zpracování dat,
			\item data se přenesou do uživatelského rozhraní,
			\item zařízení dále komunikují s ostatními prvky sítě.
		\end{itemize}
	\end{itemize}
	
	\item \textbf{M2M} (patří pod IoT)
	\begin{itemize}
		\item Komunikace mezi stroji bez zásahu člověka.
		\item Slouží k automatické výměně dat mezi zařízeními.
		\item Může probíhat bezdrátově i kabelově.
		\item Nemusí být nutně připojena k internetu.
		\item Využívá se pro přenos krátkých zpráv a provozních údajů.
		\item Bývá charakteristická vyšší latencí v řádu sekund.
		\item Má nízkou spotřebu energie, takže zařízení mohou fungovat dlouhou dobu na baterii.
	\end{itemize}
\end{itemize}

\subsection{Systém LoRa (\textit{Long Range})}
LPWAN technologie zaměřená na dlouhý dosah a nízkou spotřebu energie.\\
\begin{itemize}
	\item Možnost zapojení vlastní soukromé sítě bez poplatků operátorovi
	\item Využití rozprostření spektra (DS-SS)
	\item Přenos dat z velkého množství zařízení
	\item Modulace: CSS (Chirp Spread Spectrum) - rozprostřené spektrum s chirp signály
	\begin{itemize}	
		\item informace se nepřenáší jako jedna stálá nosná, ale jako chirp
		\item frekvence se během symbolu plynule zvyšuje nebo snižuje napříč šířkou pásma
	\end{itemize}
	\item Šířka pásma: bezlicenční 868 MHz (Evropa)
	\item Šířka kanálu: 125 kHz, 250 kHz, 500 kHz
	\item Velikost zprávy: 242 bytů
	\item např. Chytrá města, průmyslové senzory
\end{itemize}

\subsection{Systém SigFox}
\begin{itemize}
	\item LPWAN technologie určená pro velmi úsporný přenos malých zpráv na velké vzdálenosti.
	\item Funguje jako klasický operátor: uživatel využívá hotové pokrytí a nemůže si postavit vlastní základnovou stanici
	\item RF pásmo: bezlicenční 868 MHz (Evropa)
	\item Šířka kanálu: 100 Hz
	\item Modulace: GFSK (DL), DBPSK (UL)
	\item Použití: Občasné odesílání malých dat, ne pro trvalou komunikaci nebo větší objemy dat.
	\item \textbf{Architektura:}
	\begin{itemize}	
		\item využití buňkové struktury se základnovými stanicemi
		\item po přijetí zprávy ze zákaznického zařízení (modemu) zasílá základnová stanice zprávu do SigFox cloudu
		\item zprávy se dále zasílají k cílovému zařízení
	\end{itemize}
	\item Přenosový rámec tvoří jen 26 bytů (0 až 12 bytů tvoří uživatelská data; zpětný kanál omezen na 0 až 8 bytů uživatelských dat)
\end{itemize}

\subsection{LoRa vs SigFox}
SigFox je více omezený a určený hlavně pro velmi malé zprávy a minimální provoz.\\
LoRa je flexibilnější, protože umožňuje lepší kontrolu parametrů přenosu a obvykle i širší možnosti využití.\\
Obě technologie ale spadají mezi LPWAN a sdílejí hlavní cíle: dlouhý dosah, nízkou spotřebu a malý datový tok.


